% Generated draft: requires mathematical and source review.
\documentclass[11pt]{article}
\usepackage[margin=25mm]{geometry}
\usepackage{fontspec}
\setmainfont{Latin Modern Roman}
\usepackage{amsmath,amssymb,mathtools,unicode-math}
\setmathfont{Latin Modern Math}
\usepackage{enumitem,needspace,xurl,hyperref}
\setlength{\parindent}{0pt}
\setlength{\parskip}{5pt}
\setlength{\emergencystretch}{3em}
\newcommand{\N}{\mathbb{N}}
\newcommand{\Z}{\mathbb{Z}}
\newcommand{\Q}{\mathbb{Q}}
\newcommand{\R}{\mathbb{R}}
\newcommand{\C}{\mathbb{C}}
\newcommand{\F}{\mathbb{F}}
\newcommand{\A}{\mathbb{A}}
\newcommand{\PP}{\mathbb{P}}
\newcommand{\E}{\mathbb{E}}
\newcommand{\eps}{\varepsilon}
\newcommand{\dd}{\,\mathrm{d}}
\begin{document}
\section*{UnsolvedMath ID: 6200049}
\textbf{Source problem code:} catalogue entry
\subsection{Definitions and mathematical statement}
Let \(H_1,H_2\) be Hilbert spaces with their norm metrics and let \(f:H_1\to H_2\) be a homeomorphism. The source asks whether the Euclidean implication chain
\[
 \text{quasiconformal}\ \Longrightarrow\ \text{quasisymmetric}
 \ \Longrightarrow\ \text{maps balls to quasiballs}
\]
continues to hold, with the intended issue being infinite-dimensional spaces. The excerpt does not specify real versus complex Hilbert spaces, separability, dimension restrictions, or the definitions of the first and last properties.

A usual metric definition of quasiconformality uses
\[
 H_f(x)=\limsup_{r\downarrow0}
 \frac{\sup_{\|a-x\|=r}\|f(a)-f(x)\|}
 {\inf_{\|b-x\|=r}\|f(b)-f(x)\|}
 \le K\quad\text{for every }x,
\]
for a uniform finite \(K\), with suitable conventions if an extremum or denominator degenerates. This is one candidate: analytic, modulus-based, and other infinite-dimensional notions must not be silently identified with it.

A homeomorphism is \(\eta\)-quasisymmetric when an increasing homeomorphism \(\eta:[0,\infty)\to[0,\infty)\) satisfies
\[
 \frac{\|f(x)-f(a)\|}{\|f(x)-f(b)\|}
 \le\eta\left(\frac{\|x-a\|}{\|x-b\|}\right)
\]
for every triple of distinct points. The question must specify whether this is global, as written, or local, and whether control is required to depend only on \(K\).

The meaning of ``quasiball'' and uniformity over input balls are absent. A quasiconformal image of a ball, a quasisymmetric image of a ball, and a domain with controlled geometric roundness are different possible definitions; some would make parts of the proposed chain definitional. No one of them is selected here. With the intended predicates fixed, both arrows are to be tested for all eligible homeomorphisms, rather than only for a particular ball or map.

The summary labels the question open and speculates that the Euclidean chain may fail, but supplies no counterexample, theorem, or exact infinite-dimensional convention. This speculation is not treated as a negative answer, and the 2005 problem label is not an independently verified current status.
\subsection{Short English statement}
Do quasiconformal homeomorphisms of Hilbert spaces have to be quasisymmetric, and do they consequently send balls to the intended quasiballs? The infinite-dimensional definitions, global controls, and quasiball convention need to be fixed before the two implications can be stated unambiguously.
\subsection{Sources}
\begin{itemize}
\item Local registry: \url{data/unsolvedmath/problems.json}, permanent UnsolvedMath ID \(6200049\), source code catalogue entry; supplied statement, background, and research summary, updated 2026-08-06.
\item Juha Heinonen (2005), Problem 49 in Kapovich, \emph{Problems on Boundaries of Groups and Kleinian Groups}, PDF page 14: \url{https://www.math.ucdavis.edu/~kapovich/EPR/problems.pdf}. Supplied normalized PDF-text extraction accessed 2026-07-29; current status marked as needing review.
\item No separate authoritative infinite-dimensional theorem, counterexample, or quasiball definition is supplied.
\end{itemize}
\end{document}
